% ----------------------------------------------------------------------
%  3.1.5 (окончание). Ионная связь; силы Ван-дер-Ваальса; сравнение
% ----------------------------------------------------------------------
% вспомогательные макросы для рисунков этого раздела
\newcommand{\RRcell}[4]{% #1 сдвиг x, #2 полурасстояние до центра аниона, #3 радиус катиона, #4 подпись
    \begin{scope}[shift={(#1,0)}]
      \foreach \sx in {-1,1}{\foreach \sy in {-1,1}{
        \draw[atomLine, fill=atomA, line width=0.5pt] (\sx*#2,\sy*#2) circle (1);}}
      \draw[atomLine, fill=atomB, line width=0.5pt] (0,0) circle (#3);
      \node[font=\small, align=center, anchor=north] at (0,-2.45) {#4};
    \end{scope}}
\newcommand{\vdwDip}[3][1]{% [масштаб]{(центр)}{угол}: «−» слева, «+» справа
    \begin{scope}[shift={#2}, rotate=#3, scale=#1]
      \begin{scope}\clip (0,0) ellipse (0.62 and 0.3);
        \fill[fmBlue!25] (-0.7,-0.4) rectangle (0,0.4);
        \fill[fmRed!25] (0,-0.4) rectangle (0.7,0.4);
      \end{scope}
      \draw[atomLine, line width=0.5pt] (0,0) ellipse (0.62 and 0.3);
      \node[font=\scriptsize] at (-0.33,0) {$-$}; \node[font=\scriptsize] at (0.33,0) {$+$};
    \end{scope}}

\subsubsection{Ионная связь}\label{sssec:ionic}

Атомы, стоящие в периодической системе рядом с благородными газами,
стремятся приобрести их устойчивую конфигурацию. У \emph{щелочных
металлов} (сразу после благородного газа) единственный валентный
электрон движется вне заполненной оболочки, слабо связан с ядром и легко
отрывается. \emph{Галогенам} (непосредственно перед благородным газом)
до заполненной оболочки не хватает одного электрона, и они охотно его
присоединяют. Соединение таких атомов даёт \term{ионную связь}: один атом
отдаёт электрон, другой его принимает, и образовавшиеся ионы
притягиваются по закону Кулона.

\paragraph{Электроотрицательность}
Какой из двух атомов отдаст электрон, а какой примет? Это определяется их
электроотрицательностью.
\begin{opredelenie}
\term{Электроотрицательность}~— способность атома притягивать к себе
электроны при образовании химической связи. Это величина
\emph{относительная}: атом $B$ может забирать электрон у атома $A$ и
отдавать свой электрон атому $C$; тогда по электроотрицательности
$A<B<C$.
\end{opredelenie}

Мулликен предложил измерять электроотрицательность через энергию
ионизации $J$ и сродство к электрону $E$ (энергию, выделяющуюся при
присоединении электрона к нейтральному атому). Рассуждение простое.
\begin{shagi}
\item Электрон переходит от $A$ к $B$: тратится $J_A$, выделяется $E_B$;
выигрыш энергии $E_B-J_A$.
\item Обратный переход ($B\to A$) даёт выигрыш $E_A-J_B$.
\item Реализуется переход с б\'ольшим выигрышем. Переход $A\to B$ выгоднее,
если $E_B-J_A>E_A-J_B$, т.\,е. $J_B+E_B>J_A+E_A$.
\end{shagi}
Значит, электрон уходит к атому с б\'ольшей суммой $J+E$. Половину этой
суммы и называют электроотрицательностью (по Мулликену):
\begin{keyeq}\label{eq:mulliken}
\chi=\frac{J+E}{2}.
\end{keyeq}

Полинг определял электроотрицательность по энергиям диссоциации $D$
молекул $AB$, $A_2$ и $B_2$. Если связь в $AB$ чисто ковалентная, её
энергия близка к среднему из энергий $A_2$ и $B_2$; «лишняя» энергия
$\Delta=D_{AB}-\sqrt{D_{A_2}D_{B_2}}$ появляется за счёт ионного вклада и
тем больше, чем сильнее различаются атомы:
\begin{equation}\label{eq:pauling}
|\chi_A-\chi_B|=\sqrt{\Delta/(1~\text{эВ})}.
\end{equation}
Это формула для энергий в электронвольтах\ispr{В оригинале (с.~357)
коэффициент $0{,}208$ приведён вместе с размерностью «эВ». Но
$0{,}208=1/\sqrt{23{,}06}$ соответствует энергиям, выраженным в
ккал/моль ($1$~эВ на частицу${}=23{,}06$~ккал/моль); если $D$ выражено
в эВ, коэффициент равен единице.}. Всего предложено около 20 шкал
электроотрицательности, основанных на разных свойствах (межатомные
расстояния, энергии связи и~т.\,д.). Абсолютные числа в них разные, но
\emph{порядок} элементов практически одинаков, поэтому на качественные
выводы выбор шкалы не влияет. По шкале Полинга самый
электроотрицательный элемент~— фтор ($\chi\approx4{,}0$), самые
электроположительные~— цезий и франций ($\chi\approx0{,}8$).

\begin{glavnoe}
Ионная связь тем ярче выражена, чем больше разность
электроотрицательностей атомов: она образуется между атомом с большой
электроотрицательностью (галоген, кислород) и атомом с низким
потенциалом ионизации (щелочной, щелочноземельный металл). Типичные
ионные соединения~— галогениды щелочных металлов: NaCl, KBr, LiF.
\end{glavnoe}
Галогеноводороды HCl, HBr, HI, напротив, к ионным соединениям не
относятся\ispr{В оригинале (с.~357) HCl, HBr, HI названы типичными
ионными соединениями.}: разность электроотрицательностей у них невелика,
и связь в этих молекулах ковалентная полярная (доля ионности в HCl около
20\,\%, в HI~— около 5\,\%).

\paragraph{Ковалентная или ионная? Проверка на молекуле NaCl}
С классической точки зрения молекула NaCl состоит из ионов Na$^+$ и
Cl$^-$ с электронными оболочками неона и аргона, которые притягиваются
электростатически. Но так ли это? Ковалентная связь здесь тоже мыслима:
у натрия ($1s^22s^22p^63s^1$) и у хлора ($\ldots3s^23p^5$) есть по одному
неспаренному электрону, а молекулы Na$_2$, K$_2$, Li$_2$ (с ковалентной
связью, как в H$_2$) действительно существуют. Проверим обе гипотезы по
энергии связи.

\emph{Гипотеза 1: связь ковалентная.} Для ковалентной молекулы $AB$
энергия связи близка к среднему арифметическому энергий молекул $A_2$ и
$B_2$:
\begin{equation}\label{eq:cov-mean}
D_{AB}\approx\frac{D_{AA}+D_{BB}}2 .
\end{equation}
Энергии связи двухатомных молекул (эВ): Cl$_2$~— $2{,}48$; Br$_2$~—
$1{,}96$; I$_2$~— $1{,}52$; Na$_2$~— $0{,}74$; K$_2$~— $0{,}52$. Для
молекул, где связь заведомо ковалентная, правило \eqref{eq:cov-mean}
работает хорошо (табл.~\ref{tab:cov-check}). А для NaCl оно даёт
$(2{,}48+0{,}74)/2=1{,}61$~эВ при опытном значении
$4{,}26$~эВ~— в два с половиной раза меньше. Ковалентная гипотеза
отпадает: обобществления электронной пары между натрием и хлором нет.

\begin{table}[ht]
\centering
\caption{Проверка правила \eqref{eq:cov-mean}: энергии связи, эВ}
\label{tab:cov-check}
\small
\begin{tabular}{@{}lccccc@{}}
\toprule
 & BrCl & ICl & IBr & NaK & \textbf{NaCl} \\
\midrule
расчёт по \eqref{eq:cov-mean} & 2,22 & 2,00 & 1,74 & 0,63 & \textbf{1,61} \\
эксперимент & 2,26 & 2,13 & 1,83 & 0,63 & \textbf{4,26} \\
\bottomrule
\end{tabular}
\end{table}

\emph{Гипотеза 2: связь ионная.} Тогда энергию образования молекулы из
нейтральных атомов можно разбить на два этапа:
\begin{equation}\label{eq:E-ionic}
E_\text{NaCl}=E_\text{ио}+E_\text{вз},
\end{equation}
где $E_\text{ио}$~— энергия образования пары ионов из атомов, а
$E_\text{вз}$~— энергия кулоновского притяжения ионов.

Для $E_\text{ио}$ нужны энергия ионизации и \term{сродство к электрону}~—
энергия, которая выделяется ($E>0$) или затрачивается ($E<0$) при
присоединении электрона к нейтральному атому. По смыслу сродство
противоположно энергии ионизации. Электрон легко присоединяется к атомам
с почти заполненной $p$-оболочкой, особенно если при этом получается
конфигурация благородного газа (F, Cl). Атомам с полностью заполненными
оболочками (Be, Ne, Ar) пришлось бы помещать электрон на следующую,
высоколежащую оболочку~— устойчивый отрицательный ион у них не
образуется (табл.~\ref{tab:affinity}).

\begin{table}[ht]
\centering
\caption{Сродство атомов к электрону (современные данные)}
\label{tab:affinity}
\small
\begin{tabular}{@{}lccccccccc@{}}
\toprule
Элемент & H & Li & Be & F & Ne & Na & Cl & Ar & Cu \\
\midrule
$E$, эВ & 0,75 & 0,62 & $<0$ & 3,40 & $<0$ & 0,55 & 3,61 & $<0$ & 1,24 \\
\bottomrule
\end{tabular}
\end{table}

Мыслимы два варианта перераспределения электрона
(рис.~\ref{fig:NaCl-energy}):
\begin{shagi}
\item Na$^+$ + Cl$^-$: $E_\text{ио}=-J_\text{Na}+E_\text{Cl}=-5{,}14+3{,}61=-1{,}53$~эВ;
\item Na$^-$ + Cl$^+$: $E_\text{ио}=E_\text{Na}-J_\text{Cl}=0{,}55-12{,}97=-12{,}42$~эВ.
\end{shagi}
Оба варианта сами по себе невыгодны (энергию надо затратить), но первый
обходится в восемь раз «дешевле». Кулоновский выигрыш одинаков для обоих
вариантов; при равновесном расстоянии в молекуле $r_0=0{,}236$~нм
\[
E_\text{вз}=\frac{e^2}{4\pi\varepsilon_0r_0}
=\frac{1{,}44~\text{эВ}\cdot\text{нм}}{0{,}236~\text{нм}}\approx6{,}10~\text{эВ}.
\]
Он с запасом покрывает затраты в первом варианте, но не во втором. Итак,
энергия связи ионной молекулы NaCl
\[
E_\text{NaCl}\approx6{,}10-1{,}53=4{,}57~\text{эВ},
\]
а эксперимент даёт $4{,}26$~эВ\ispr{В оригинале использованы устаревшие
данные: $J_\text{Na}=5{,}12$~эВ, $E_\text{Cl}=3{,}7$~эВ,
$E_\text{Na}=0{,}22$~эВ, $E_\text{вз}=5{,}78$~эВ (это соответствует
$r_0\approx0{,}25$~нм), и получено $E_\text{NaCl}=4{,}36$~эВ. Кроме того,
во втором варианте допущена арифметическая ошибка: $0{,}22-13{,}01=-12{,}79$,
а не $-13{,}82$~эВ. На выводы это не влияет.}. Согласие хорошее (около
7\,\%), а небольшой избыток расчётной энергии объясняется тем, что мы
не учли отталкивание электронных оболочек ионов на малых расстояниях
(п.~\ref{ssec:binding}).

\begin{figure}[ht]
\centering
\begin{tikzpicture}[yscale=0.52, xscale=1.0,
  lev/.style={line width=1.4pt}, st/.style={-{Stealth[length=2.2mm]}, line width=0.8pt}]
  \draw[-{Stealth[length=2mm]}, fmGray] (-0.4,-5.4) -- (-0.4,6.2) node[above, font=\small]{$E$, эВ};
  \foreach \y in {-4,-2,0,2,4}{\draw[fmGray] (-0.48,\y) -- (-0.32,\y) node[left=2pt, font=\scriptsize]{$\y$};}
  % уровни
  \draw[lev] (0,0) -- (2.3,0) node[midway, above, font=\small]{Na + Cl};
  \draw[lev, fmRed] (2.7,5.14) -- (5.0,5.14) node[midway, above, font=\small]{Na$^+$ + $e^-$ + Cl};
  \draw[lev, fmAmber] (5.4,1.53) -- (7.7,1.53) node[midway, above, font=\small, align=center]{Na$^+$ + Cl$^-$\\[-2pt]\scriptsize (ионы далеко)};
  \draw[lev, fmBlue] (8.1,-4.57) -- (10.4,-4.57) node[midway, below, font=\small]{NaCl (расчёт)};
  \draw[lev, fmBlue, dashed, line width=0.9pt] (8.1,-4.26) -- (10.4,-4.26) node[midway, above, font=\scriptsize]{эксперимент};
  % переходы
  \draw[st, fmRed] (2.5,0) -- (2.5,5.14) node[midway, left, font=\scriptsize, align=right]{ионизация\\Na: $+5{,}14$};
  \draw[st, fmAmber] (5.2,5.14) -- (5.2,1.53) node[midway, left, font=\scriptsize, align=right]{сродство\\Cl: $-3{,}61$};
  \draw[st, fmBlue] (7.9,1.53) -- (7.9,-4.57) node[midway, left, font=\scriptsize, align=right]{кулоновское\\притяжение:\\$-6{,}10$};
  \draw[dashed, fmGray] (2.3,0) -- (10.8,0);
  \draw[dimline] (10.8,0) -- (10.8,-4.57) node[midway, right, font=\scriptsize, align=left]{энергия\\связи\\$4{,}57$};
\end{tikzpicture}
\caption{Энергетический баланс образования молекулы NaCl из атомов.
Образование пары ионов требует затраты $5{,}14-3{,}61=1{,}53$~эВ, но
кулоновское притяжение ионов возвращает $6{,}10$~эВ}
\label{fig:NaCl-energy}
\end{figure}

\begin{glavnoe}
Ионная связь~— результат \emph{перехода} валентного электрона от
электроположительного атома к электроотрицательному и последующего
кулоновского притяжения ионов. Ионы имеют сферически симметричные
оболочки благородных газов, поэтому ионная связь, как и металлическая,
\vazhno{ненаправленна и ненасыщаема}.
\end{glavnoe}

\paragraph{Строение ионных кристаллов}
Анион получается присоединением электронов к нейтральному атому, катион~—
их отрывом. Поэтому анионы, как правило, заметно крупнее катионов. Отсюда
простая картина ионного кристалла:
\begin{itemize}
\item основной объём заполняют крупные \emph{анионы}; связь ненаправленна
и ненасыщаема, поэтому они образуют плотную шаровую упаковку;
\item мелкие \emph{катионы} располагаются в пустотах этой упаковки,
подходящих им по размеру;
\item ближайшие соседи каждого иона~— ионы противоположного знака.
\end{itemize}
Чтобы понять, какие пустоты займут катионы, нужно знать размеры пустот.

\paragraph{Пустоты в плотных упаковках}
В плотнейших упаковках шаров (ГЦК и ГПУ; см. п.~1.4) есть пустоты трёх
типов: треугольные, тетраэдрические и октаэдрические. Для ионных структур
понадобится ещё кубическая пустота~— в центре простой кубической ячейки
(рис.~\ref{fig:voids}). \term{Радиусом пустоты} называют радиус шара,
который можно в неё поместить, не раздвигая окружающих шаров.

\begin{figure}[ht]
\centering
\begin{subfigure}[b]{0.22\textwidth}\centering
\begin{tikzpicture}[scale=0.62]
  \foreach \a in {90,210,330}{\draw[atomLine, fill=atomA, line width=0.5pt] (\a:1.1547) circle (1);}
  \draw[atomLine, fill=atomB, line width=0.5pt] (0,0) circle (0.1547);
  \draw[fmRed, line width=0.5pt] (90:1.1547) -- (210:1.1547) -- (330:1.1547) -- cycle;
\end{tikzpicture}
\caption{$r/R=0{,}155$}
\end{subfigure}\hfill
\begin{subfigure}[b]{0.24\textwidth}\centering
\begin{tikzpicture}[cubeproj, scale=1.9]
  \draw[atomLine!30, line width=0.4pt] (0,0,0)--(1,0,0)--(1,0,1)--(0,0,1)--cycle (1,0,0)--(1,1,0)--(1,1,1)--(1,0,1)
     (0,0,1)--(0,1,1)--(1,1,1);
  \draw[fmRed, line width=0.6pt] (0,0,0)--(1,1,0)--(1,0,1)--cycle (0,1,1)--(1,1,0) (0,1,1)--(1,0,1);
  \draw[fmRed, line width=0.6pt, dashed] (0,0,0)--(0,1,1);
  \node[atom=4mm] at (0,1,1) {};
  \node[atom=4mm] at (1,1,0) {};
  \node[atomB=2.4mm] at (0.5,0.5,0.5) {};
  \node[atom=4mm] at (0,0,0) {};
  \node[atom=4mm] at (1,0,1) {};
\end{tikzpicture}
\caption{$r/R=0{,}225$}
\end{subfigure}\hfill
\begin{subfigure}[b]{0.24\textwidth}\centering
\begin{tikzpicture}[cubeproj, scale=1.9]
  \draw[fmRed, line width=0.6pt, dashed] (0.5,1,0.5) -- (0,0.5,0.5) (0.5,1,0.5) -- (1,0.5,0.5)
     (0.5,1,0.5) -- (0.5,0.5,0) (0.5,1,0.5) -- (0.5,0.5,1);
  \draw[fmRed, line width=0.6pt] (0.5,0,0.5) -- (0,0.5,0.5) (0.5,0,0.5) -- (1,0.5,0.5)
     (0.5,0,0.5) -- (0.5,0.5,0) (0.5,0,0.5) -- (0.5,0.5,1)
     (0,0.5,0.5) -- (0.5,0.5,1) -- (1,0.5,0.5) -- (0.5,0.5,0) -- cycle;
  \node[atom=4mm] at (0.5,1,0.5) {};
  \node[atom=4mm] at (0,0.5,0.5) {};
  \node[atom=4mm] at (0.5,0.5,0) {};
  \node[atomB=2.4mm] at (0.5,0.5,0.5) {};
  \node[atom=4mm] at (1,0.5,0.5) {};
  \node[atom=4mm] at (0.5,0.5,1) {};
  \node[atom=4mm] at (0.5,0,0.5) {};
\end{tikzpicture}
\caption{$r/R=0{,}414$}
\end{subfigure}\hfill
\begin{subfigure}[b]{0.24\textwidth}\centering
\begin{tikzpicture}[cubeproj, scale=1.6]
  \CubeEdges[draw=fmRed]
  \foreach \p in {(0,1,0),(1,1,0),(0,1,1),(1,1,1)} \node[atom=4mm] at \p {};
  \node[atomB=2.4mm] at (0.5,0.5,0.5) {};
  \foreach \p in {(0,0,0),(1,0,0),(0,0,1),(1,0,1)} \node[atom=4mm] at \p {};
\end{tikzpicture}
\caption{$r/R=0{,}732$}
\end{subfigure}
\caption{Пустоты в упаковках шаров: \textit{а}~— треугольная (точно в
масштабе), \textit{б}~— тетраэдрическая, \textit{в}~— октаэдрическая,
\textit{г}~— кубическая. Голубые шары~— атомы упаковки (их размеры на
объёмных рисунках уменьшены), оранжевый~— шар, помещающийся в пустоту}
\label{fig:voids}
\end{figure}

Радиусы пустот находятся из простой геометрии. Примем радиус шаров
упаковки $R=1$; тогда расстояние между центрами соседних шаров равно~2,
и радиус пустоты $r$ равен расстоянию от её центра до центра шара минус
единица.

\begin{primer}[Радиус октаэдрической пустоты]
Октаэдр образован шестью шарами; в его сечении через центр~— квадрат из
четырёх шаров со стороной $2R=2$ (рис.~\ref{fig:radius-ratio},\,\textit{б}).
Расстояние от центра квадрата до вершины~— половина диагонали:
$\frac12\cdot2\sqrt2=\sqrt2$. Значит,
\[
r_\text{окт}=\sqrt2-1\approx0{,}414 .
\]
\end{primer}

Аналогично для остальных пустот (табл.~\ref{tab:voids}): в треугольнике
со стороной~2 расстояние от центра до вершины равно $2/\sqrt3$; в
тетраэдре с ребром~2~— $\sqrt{3/2}$; в кубе с ребром~2~— половина
пространственной диагонали, $\sqrt3$.

\begin{table}[ht]
\centering
\caption{Радиусы пустот при радиусе шаров упаковки $R=1$}
\label{tab:voids}
\small
\begin{tabular}{@{}lccl@{}}
\toprule
Пустота & Соседей & Радиус $r$ & Где встречается \\
\midrule
треугольная & 3 & $2/\sqrt3-1=0{,}155$ & в слое плотной упаковки \\
тетраэдрическая & 4 & $\sqrt{3/2}-1=0{,}225$ & ГЦК, ГПУ: 2 на атом \\
октаэдрическая & 6 & $\sqrt2-1=0{,}414$ & ГЦК, ГПУ: 1 на атом \\
кубическая & 8 & $\sqrt3-1=0{,}732$ & простая кубическая: 1~на~атом \\
\bottomrule
\end{tabular}
\end{table}

Расположение тетраэдрических и октаэдрических пустот в ГЦК- и
ГПУ-решётках показано в главе~1 (п.~1.4). Важно запомнить: в плотнейших
упаковках \vazhno{на каждый атом приходятся одна октаэдрическая и две
тетраэдрические пустоты}.

\paragraph{Правило отношения радиусов}
Пусть небольшой катион окружён анионами. Сколько анионов он может
удержать вокруг себя, зависит от отношения радиусов катиона и аниона
$r_\text{к}/r_\text{а}$ (рис.~\ref{fig:radius-ratio}).
\begin{itemize}
\item Катион \emph{больше} пустоты: он касается всех анионов и слегка
раздвигает их. Разноимённые ионы в контакте~— конфигурация устойчива.
\item Катион \emph{точно} равен пустоте: касаются все ионы~— это
предельный случай.
\item Катион \emph{меньше} пустоты: анионы касаются друг друга, а катион
«болтается» в пустоте и не касается их. Разноимённые заряды оказываются
дальше, чем могли бы быть, и энергия выше. Выгоднее перейти к меньшему
координационному числу, при котором анионы снова прижмутся к катиону.
\end{itemize}
Обратно: если катион настолько велик, что вокруг него помещается больше
анионов, структура перестраивается с увеличением координационного числа.
Отсюда предельные значения отношения радиусов
(табл.~\ref{tab:radius-ratio})~— они в точности совпадают с радиусами
пустот из табл.~\ref{tab:voids}.

\begin{figure}[ht]
\centering
\begin{tikzpicture}[scale=0.62]
  \RRcell{0}{1.1314}{0.6}{\textit{а}~$r_\text{к}/r_\text{а}=0{,}6$:\\устойчиво}
  \RRcell{6.2}{1.0}{0.414}{\textit{б}~$r_\text{к}/r_\text{а}=0{,}414$:\\предел}
  \RRcell{12.4}{1.0}{0.25}{\textit{в}~$r_\text{к}/r_\text{а}=0{,}25$:\\неустойчиво}
  \draw[fmRed, line width=0.5pt] (12.4,0) circle (0.414);
\end{tikzpicture}
\caption{Сечение октаэдрической координации (четыре аниона в плоскости;
ещё два~— над и под катионом). В случае \textit{в} красная окружность
показывает размер пустоты: катион в ней не касается анионов, и выгоднее
тетраэдрическая координация}
\label{fig:radius-ratio}
\end{figure}

\begin{table}[ht]
\centering
\caption{Предельные отношения ионных радиусов для симметричных координаций}
\label{tab:radius-ratio}
\small
\begin{tabular}{@{}cll@{}}
\toprule
Координационное число & Координация & $r_\text{к}/r_\text{а}$ \\
\midrule
3 & треугольник & 0,155--0,225 \\
4 & тетраэдр & 0,225--0,414 \\
6 & октаэдр & 0,414--0,732 \\
8 & куб & 0,732--1 \\
\bottomrule
\end{tabular}
\end{table}

Ещё одно ограничение накладывает \emph{электронейтральность}: суммарные
заряды катионов и анионов должны быть равны. Поэтому пустоты нужного
типа часто заполняются не все, а лишь их часть (половина, четверть и~т.\,д.).

\paragraph{Важнейшие структурные типы ионных кристаллов}
Посмотрим, как эти правила работают на примерах (рис.~\ref{fig:ionic-cells},
табл.~\ref{tab:ionic-types}).

\begin{itemize}
\item \term{Структура NaCl}~— ГЦК-упаковка анионов Cl$^-$, все
октаэдрические пустоты заняты катионами Na$^+$. Вдоль всех трёх осей куба
ионы натрия и хлора чередуются. Радиусы: $r_{\text{Na}^+}=0{,}098$~нм,
$r_{\text{Cl}^-}=0{,}181$~нм, отношение $0{,}54$~— как раз в интервале
октаэдрической координации. Радиус октаэдрической пустоты в упаковке
ионов хлора $0{,}414\cdot0{,}181\approx0{,}075$~нм меньше радиуса Na$^+$,
поэтому катионы слегка раздвигают анионы (случай \textit{а} на
рис.~\ref{fig:radius-ratio}).
\item \term{Структура NiAs}~— то же, но на основе ГПУ-упаковки анионов:
катионы занимают все октаэдрические пустоты. Как и в NaCl, на каждый
анион приходится один катион.
\item \term{Сфалерит} ($\alpha$-ZnS, цинковая обманка)~— ГЦК-упаковка
анионов S$^{2-}$, катионы Zn$^{2+}$ занимают \emph{половину}
тетраэдрических пустот (центры четырёх из восьми малых кубиков).
Расположение то же, что в алмазе: если заменить все атомы цинка и серы
на углерод, получится решётка алмаза (рис.~\ref{fig:diamond}). Поэтому
сфалерит называют \emph{алмазоподобной} структурой.
\item \term{Вюрцит} ($\beta$-ZnS, гексагональная модификация)~—
ГПУ-упаковка анионов S$^{2-}$, катионы занимают половину тетраэдрических
пустот.
\item \term{Флюорит} (CaF$_2$, плавиковый шпат\ispr{В оригинале (с.~364)
ошибочно: «кольцевой шпат».}): $r_{\text{Ca}^{2+}}=0{,}104$~нм,
$r_{\text{F}^-}=0{,}133$~нм, отношение $0{,}78$~— должны заполняться
кубические пустоты. Анионы F$^-$ образуют простую кубическую решётку,
катионы занимают \emph{половину} кубических пустот (электронейтральность:
фторов вдвое больше). Ту же структуру можно описать иначе: катионы
Ca$^{2+}$ образуют ГЦК-решётку, а анионы F$^-$ сидят в центрах всех
восьми малых кубиков (рис.~\ref{fig:ionic-cells},\,\textit{г}). Такую
структуру имеет и главное ядерное топливо~— диоксид урана UO$_2$.
\item \term{Структура CsCl}: $r_{\text{Cs}^+}=0{,}165$~нм,
$r_{\text{Cl}^-}=0{,}181$~нм, отношение $0{,}91$~— кубическая
координация. Анионы образуют простую кубическую решётку, катионы
занимают \emph{все} кубические пустоты.
\end{itemize}

\begin{figure}[ht]
\centering
\begin{subfigure}[b]{0.24\textwidth}\centering
\begin{tikzpicture}[cubeproj, scale=1.85, anion/.style={atom=4.4mm}, cation/.style={atomB=2.8mm}]
  \CubeEdges[draw opacity=0]
  % сгенерировано tools/gen_ionic_cells.py — не править вручную
\draw[bondline] (0,0,0) -- (0,0,0.5);
\draw[bondline] (0,0,0) -- (0,0.5,0);
\draw[bondline] (0,0,0) -- (0.5,0,0);
\draw[bondline] (0,0,0.5) -- (0,0,1);
\draw[bondline] (0,0,0.5) -- (0,0.5,0.5);
\draw[bondline] (0,0,0.5) -- (0.5,0,0.5);
\draw[bondline] (0,0,1) -- (0,0.5,1);
\draw[bondline] (0,0,1) -- (0.5,0,1);
\draw[bondline] (0,0.5,0) -- (0,0.5,0.5);
\draw[bondline] (0,0.5,0) -- (0,1,0);
\draw[bondline] (0,0.5,0) -- (0.5,0.5,0);
\draw[bondline] (0,0.5,0.5) -- (0,0.5,1);
\draw[bondline] (0,0.5,0.5) -- (0,1,0.5);
\draw[bondline] (0,0.5,0.5) -- (0.5,0.5,0.5);
\draw[bondline] (0,0.5,1) -- (0,1,1);
\draw[bondline] (0,0.5,1) -- (0.5,0.5,1);
\draw[bondline] (0,1,0) -- (0,1,0.5);
\draw[bondline] (0,1,0) -- (0.5,1,0);
\draw[bondline] (0,1,0.5) -- (0,1,1);
\draw[bondline] (0,1,0.5) -- (0.5,1,0.5);
\draw[bondline] (0,1,1) -- (0.5,1,1);
\draw[bondline] (0.5,0,0) -- (0.5,0,0.5);
\draw[bondline] (0.5,0,0) -- (0.5,0.5,0);
\draw[bondline] (0.5,0,0) -- (1,0,0);
\draw[bondline] (0.5,0,0.5) -- (0.5,0,1);
\draw[bondline] (0.5,0,0.5) -- (0.5,0.5,0.5);
\draw[bondline] (0.5,0,0.5) -- (1,0,0.5);
\draw[bondline] (0.5,0,1) -- (0.5,0.5,1);
\draw[bondline] (0.5,0,1) -- (1,0,1);
\draw[bondline] (0.5,0.5,0) -- (0.5,0.5,0.5);
\draw[bondline] (0.5,0.5,0) -- (0.5,1,0);
\draw[bondline] (0.5,0.5,0) -- (1,0.5,0);
\draw[bondline] (0.5,0.5,0.5) -- (0.5,0.5,1);
\draw[bondline] (0.5,0.5,0.5) -- (0.5,1,0.5);
\draw[bondline] (0.5,0.5,0.5) -- (1,0.5,0.5);
\draw[bondline] (0.5,0.5,1) -- (0.5,1,1);
\draw[bondline] (0.5,0.5,1) -- (1,0.5,1);
\draw[bondline] (0.5,1,0) -- (0.5,1,0.5);
\draw[bondline] (0.5,1,0) -- (1,1,0);
\draw[bondline] (0.5,1,0.5) -- (0.5,1,1);
\draw[bondline] (0.5,1,0.5) -- (1,1,0.5);
\draw[bondline] (0.5,1,1) -- (1,1,1);
\draw[bondline] (1,0,0) -- (1,0,0.5);
\draw[bondline] (1,0,0) -- (1,0.5,0);
\draw[bondline] (1,0,0.5) -- (1,0,1);
\draw[bondline] (1,0,0.5) -- (1,0.5,0.5);
\draw[bondline] (1,0,1) -- (1,0.5,1);
\draw[bondline] (1,0.5,0) -- (1,0.5,0.5);
\draw[bondline] (1,0.5,0) -- (1,1,0);
\draw[bondline] (1,0.5,0.5) -- (1,0.5,1);
\draw[bondline] (1,0.5,0.5) -- (1,1,0.5);
\draw[bondline] (1,0.5,1) -- (1,1,1);
\draw[bondline] (1,1,0) -- (1,1,0.5);
\draw[bondline] (1,1,0.5) -- (1,1,1);
\node[anion] at (0,1,0) {};
\node[cation] at (0.5,1,0) {};
\node[anion] at (1,1,0) {};
\node[cation] at (0,1,0.5) {};
\node[anion] at (0.5,1,0.5) {};
\node[cation] at (1,1,0.5) {};
\node[anion] at (0,1,1) {};
\node[cation] at (0.5,1,1) {};
\node[anion] at (1,1,1) {};
\node[cation] at (0,0.5,0) {};
\node[anion] at (0.5,0.5,0) {};
\node[cation] at (1,0.5,0) {};
\node[anion] at (0,0.5,0.5) {};
\node[cation] at (0.5,0.5,0.5) {};
\node[anion] at (1,0.5,0.5) {};
\node[cation] at (0,0.5,1) {};
\node[anion] at (0.5,0.5,1) {};
\node[cation] at (1,0.5,1) {};
\node[anion] at (0,0,0) {};
\node[cation] at (0.5,0,0) {};
\node[anion] at (1,0,0) {};
\node[cation] at (0,0,0.5) {};
\node[anion] at (0.5,0,0.5) {};
\node[cation] at (1,0,0.5) {};
\node[anion] at (0,0,1) {};
\node[cation] at (0.5,0,1) {};
\node[anion] at (1,0,1) {};

\end{tikzpicture}
\caption{NaCl}
\end{subfigure}\hfill
\begin{subfigure}[b]{0.24\textwidth}\centering
\begin{tikzpicture}[cubeproj, scale=1.6, anion/.style={atom=5mm}, cation/.style={atomB=4.6mm}]
  \CubeEdges
  % сгенерировано tools/gen_ionic_cells.py — не править вручную
\draw[bondline] (0.5,0.5,0.5) -- (0,0,0);
\draw[bondline] (0.5,0.5,0.5) -- (0,0,1);
\draw[bondline] (0.5,0.5,0.5) -- (0,1,0);
\draw[bondline] (0.5,0.5,0.5) -- (0,1,1);
\draw[bondline] (0.5,0.5,0.5) -- (1,0,0);
\draw[bondline] (0.5,0.5,0.5) -- (1,0,1);
\draw[bondline] (0.5,0.5,0.5) -- (1,1,0);
\draw[bondline] (0.5,0.5,0.5) -- (1,1,1);
\node[anion] at (0,1,0) {};
\node[anion] at (1,1,0) {};
\node[anion] at (0,1,1) {};
\node[anion] at (1,1,1) {};
\node[cation] at (0.5,0.5,0.5) {};
\node[anion] at (0,0,0) {};
\node[anion] at (1,0,0) {};
\node[anion] at (0,0,1) {};
\node[anion] at (1,0,1) {};

\end{tikzpicture}
\caption{CsCl}
\end{subfigure}\hfill
\begin{subfigure}[b]{0.24\textwidth}\centering
\begin{tikzpicture}[cubeproj, scale=1.85, anion/.style={atom=4.2mm}, cation/.style={atomB=2.8mm}]
  \CubeEdges
  % сгенерировано tools/gen_ionic_cells.py — не править вручную
\draw[bondthick] (0.25,0.25,0.25) -- (0,0,0);
\draw[bondthick] (0.25,0.25,0.25) -- (0,0.5,0.5);
\draw[bondthick] (0.25,0.25,0.25) -- (0.5,0,0.5);
\draw[bondthick] (0.25,0.25,0.25) -- (0.5,0.5,0);
\draw[bondthick] (0.75,0.75,0.25) -- (0.5,0.5,0);
\draw[bondthick] (0.75,0.75,0.25) -- (0.5,1,0.5);
\draw[bondthick] (0.75,0.75,0.25) -- (1,0.5,0.5);
\draw[bondthick] (0.75,0.75,0.25) -- (1,1,0);
\draw[bondthick] (0.75,0.25,0.75) -- (0.5,0,0.5);
\draw[bondthick] (0.75,0.25,0.75) -- (0.5,0.5,1);
\draw[bondthick] (0.75,0.25,0.75) -- (1,0,1);
\draw[bondthick] (0.75,0.25,0.75) -- (1,0.5,0.5);
\draw[bondthick] (0.25,0.75,0.75) -- (0,0.5,0.5);
\draw[bondthick] (0.25,0.75,0.75) -- (0,1,1);
\draw[bondthick] (0.25,0.75,0.75) -- (0.5,0.5,1);
\draw[bondthick] (0.25,0.75,0.75) -- (0.5,1,0.5);
\node[anion] at (0,1,0) {};
\node[anion] at (1,1,0) {};
\node[anion] at (0.5,1,0.5) {};
\node[anion] at (0,1,1) {};
\node[anion] at (1,1,1) {};
\node[cation] at (0.75,0.75,0.25) {};
\node[cation] at (0.25,0.75,0.75) {};
\node[anion] at (0.5,0.5,0) {};
\node[anion] at (0,0.5,0.5) {};
\node[anion] at (1,0.5,0.5) {};
\node[anion] at (0.5,0.5,1) {};
\node[cation] at (0.25,0.25,0.25) {};
\node[cation] at (0.75,0.25,0.75) {};
\node[anion] at (0,0,0) {};
\node[anion] at (1,0,0) {};
\node[anion] at (0.5,0,0.5) {};
\node[anion] at (0,0,1) {};
\node[anion] at (1,0,1) {};

\end{tikzpicture}
\caption{ZnS (сфалерит)}
\end{subfigure}\hfill
\begin{subfigure}[b]{0.24\textwidth}\centering
\begin{tikzpicture}[cubeproj, scale=1.85, anion/.style={atom=3.8mm}, cation/.style={atomB=3.4mm}]
  \CubeEdges
  % сгенерировано tools/gen_ionic_cells.py — не править вручную
\draw[bondline] (0.25,0.25,0.25) -- (0,0,0);
\draw[bondline] (0.25,0.25,0.25) -- (0,0.5,0.5);
\draw[bondline] (0.25,0.25,0.25) -- (0.5,0,0.5);
\draw[bondline] (0.25,0.25,0.25) -- (0.5,0.5,0);
\draw[bondline] (0.25,0.25,0.75) -- (0,0,1);
\draw[bondline] (0.25,0.25,0.75) -- (0,0.5,0.5);
\draw[bondline] (0.25,0.25,0.75) -- (0.5,0,0.5);
\draw[bondline] (0.25,0.25,0.75) -- (0.5,0.5,1);
\draw[bondline] (0.25,0.75,0.25) -- (0,0.5,0.5);
\draw[bondline] (0.25,0.75,0.25) -- (0,1,0);
\draw[bondline] (0.25,0.75,0.25) -- (0.5,0.5,0);
\draw[bondline] (0.25,0.75,0.25) -- (0.5,1,0.5);
\draw[bondline] (0.25,0.75,0.75) -- (0,0.5,0.5);
\draw[bondline] (0.25,0.75,0.75) -- (0,1,1);
\draw[bondline] (0.25,0.75,0.75) -- (0.5,0.5,1);
\draw[bondline] (0.25,0.75,0.75) -- (0.5,1,0.5);
\draw[bondline] (0.75,0.25,0.25) -- (0.5,0,0.5);
\draw[bondline] (0.75,0.25,0.25) -- (0.5,0.5,0);
\draw[bondline] (0.75,0.25,0.25) -- (1,0,0);
\draw[bondline] (0.75,0.25,0.25) -- (1,0.5,0.5);
\draw[bondline] (0.75,0.25,0.75) -- (0.5,0,0.5);
\draw[bondline] (0.75,0.25,0.75) -- (0.5,0.5,1);
\draw[bondline] (0.75,0.25,0.75) -- (1,0,1);
\draw[bondline] (0.75,0.25,0.75) -- (1,0.5,0.5);
\draw[bondline] (0.75,0.75,0.25) -- (0.5,0.5,0);
\draw[bondline] (0.75,0.75,0.25) -- (0.5,1,0.5);
\draw[bondline] (0.75,0.75,0.25) -- (1,0.5,0.5);
\draw[bondline] (0.75,0.75,0.25) -- (1,1,0);
\draw[bondline] (0.75,0.75,0.75) -- (0.5,0.5,1);
\draw[bondline] (0.75,0.75,0.75) -- (0.5,1,0.5);
\draw[bondline] (0.75,0.75,0.75) -- (1,0.5,0.5);
\draw[bondline] (0.75,0.75,0.75) -- (1,1,1);
\node[cation] at (0,1,0) {};
\node[cation] at (1,1,0) {};
\node[cation] at (0.5,1,0.5) {};
\node[cation] at (0,1,1) {};
\node[cation] at (1,1,1) {};
\node[anion] at (0.25,0.75,0.25) {};
\node[anion] at (0.75,0.75,0.25) {};
\node[anion] at (0.25,0.75,0.75) {};
\node[anion] at (0.75,0.75,0.75) {};
\node[cation] at (0.5,0.5,0) {};
\node[cation] at (0,0.5,0.5) {};
\node[cation] at (1,0.5,0.5) {};
\node[cation] at (0.5,0.5,1) {};
\node[anion] at (0.25,0.25,0.25) {};
\node[anion] at (0.75,0.25,0.25) {};
\node[anion] at (0.25,0.25,0.75) {};
\node[anion] at (0.75,0.25,0.75) {};
\node[cation] at (0,0,0) {};
\node[cation] at (1,0,0) {};
\node[cation] at (0.5,0,0.5) {};
\node[cation] at (0,0,1) {};
\node[cation] at (1,0,1) {};

\end{tikzpicture}
\caption{CaF$_2$ (флюорит)}
\end{subfigure}
\caption{Элементарные ячейки важнейших ионных структур. Голубые шары~—
анионы, оранжевые~— катионы (размеры уменьшены, чтобы была видна
внутренность ячейки). Линии в \textit{а}, \textit{б}, \textit{г} соединяют
ближайших соседей; в \textit{в} показаны тетраэдрические связи Zn--S}
\label{fig:ionic-cells}
\end{figure}

\begin{table}[ht]
\centering
\caption{Структурные типы ионных кристаллов}
\label{tab:ionic-types}
\small
\begin{tabularx}{\textwidth}{@{}lllc>{\raggedright\arraybackslash}X@{}}
\toprule
Тип & Упаковка анионов & Какие пустоты заняты & К.\,ч. & Примеры \\
\midrule
NaCl & ГЦК & все октаэдрические & 6 & NaCl, KCl, LiF, MgO \\
NiAs & ГПУ & все октаэдрические & 6 & NiAs, FeS \\
сфалерит & ГЦК & половина тетраэдрических & 4 & $\alpha$-ZnS, GaAs \\
вюрцит & ГПУ & половина тетраэдрических & 4 & $\beta$-ZnS, ZnO \\
флюорит & простая кубическая & половина кубических & 8 & CaF$_2$, UO$_2$ \\
CsCl & простая кубическая & все кубические & 8 & CsCl, CsBr \\
\bottomrule
\end{tabularx}
\end{table}

\begin{vnimanie}[Правило~— не закон]
Правило отношения радиусов~— геометрическая подсказка, а не строгий закон.
Оно хорошо работает для сильно ионных соединений, но может нарушаться,
если в связи заметна ковалентная доля. Пример~— сам ZnS: отношение
$r_{\text{Zn}^{2+}}/r_{\text{S}^{2-}}\approx0{,}46$ «требует»
октаэдров, а реализуется тетраэдрическая координация: направленные
$sp^3$-связи, как в алмазе, оказываются важнее упаковки шаров.
\end{vnimanie}

\paragraph{Ионные радиусы}
В ионном кристалле ионы разных знаков касаются друг друга, поэтому
межатомное расстояние равно сумме ионных радиусов:
\begin{equation}\label{eq:ionic-sum}
d_{A^+B^-}=r_{A^+}+r_{B^-} .
\end{equation}
Расстояние $d$ легко измерить дифракционными методами. Но как разделить
его на два радиуса?

Сначала убедимся, что ионные радиусы вообще имеют смысл, т.\,е. не
зависят от партнёра. Сравним соединения калия и натрия с одними и теми
же галогенами (все они имеют структуру NaCl). Разности межатомных
расстояний (нм):
\[
d_\text{KF}-d_\text{NaF}=0{,}036,\quad d_\text{KCl}-d_\text{NaCl}=0{,}033,\quad
d_\text{KBr}-d_\text{NaBr}=0{,}031,\quad d_\text{KI}-d_\text{NaI}=0{,}030 .
\]
Разность почти постоянна, в среднем $0{,}033$~нм. По
\eqref{eq:ionic-sum} это и есть $r_{\text{K}^+}-r_{\text{Na}^+}$: радиус
галогена сокращается. Вывод: \vazhno{если тип структуры не меняется,
размер иона не зависит от того, с каким ионом он соседствует}, и радиусы
можно свести в таблицу.

Чтобы получить сами радиусы, а не только их разности, нужно знать радиус
хотя бы одного иона. Так и строились таблицы ионных радиусов:
\begin{itemize}
\item \emph{Гольдшмидт} (первая таблица) взял опытные значения
$r_{\text{O}^{2-}}=0{,}132$~нм и $r_{\text{F}^-}=0{,}133$~нм;
\item \emph{Полинг}~— $r_{\text{O}^{2-}}=0{,}140$~нм и $r_{\text{F}^-}=0{,}136$~нм;
\item \emph{Г.\,Б.~Бокий и Н.\,В.~Белов} в нашей стране~— средние значения
между данными Гольдшмидта и Полинга.
\end{itemize}
Во всех таблицах радиусы приведены к координационному числу~6 и совпадают
между собой в пределах $\pm3$\,\%.

\paragraph{Изоэлектронные ряды}
У ионов O$^{2-}$, F$^-$, Na$^+$, Mg$^{2+}$, Al$^{3+}$, Si$^{4+}$ одинаковая
электронная оболочка~— как у неона (10 электронов). Тем не менее их
радиусы сильно различаются (рис.~\ref{fig:isoelectronic}): по мере роста
заряда ядра те же 10 электронов притягиваются всё сильнее, и ион
сжимается. Кроме того:
\begin{itemize}
\item в каждой группе радиусы одинаково построенных ионов растут с
атомным номером (добавляются новые оболочки);
\item из-за лантаноидного сжатия (п.~\ref{ssec:atoms}) ионы элементов
шестого периода почти равны по размеру ионам пятого.
\end{itemize}

\begin{figure}[ht]
\centering
\begin{tikzpicture}[x=1mm, y=1mm]
  \def\sc{60}% мм на нм
  \foreach \ion/\z/\r/\x/\col in {O$^{2-}$/8/0.136/0/atomA, F$^-$/9/0.133/18/atomA,
     Na$^+$/11/0.098/38/atomB, Mg$^{2+}$/12/0.074/53/atomB,
     Al$^{3+}$/13/0.057/66/atomB, Si$^{4+}$/14/0.039/78/atomB}{
    \draw[atomLine, fill=\col, line width=0.5pt] (\x,0) circle ({\r*\sc});
    \node[font=\small] at (\x,-12) {\ion};
    \node[font=\scriptsize, fmGray] at (\x,-16.5) {$Z=\z$};
    \node[font=\scriptsize, fmGray] at (\x,-20.5) {\r};}
  \node[font=\scriptsize, fmGray, anchor=east] at (-10,-20.5) {$r$, нм};
\end{tikzpicture}
\caption{Изоэлектронный ряд: у всех ионов оболочка неона (10 электронов),
но с ростом заряда ядра $Z$ радиус иона уменьшается более чем втрое
(радиусы по Бокию и Белову, в масштабе)}
\label{fig:isoelectronic}
\end{figure}

\paragraph{Изоморфизм и морфотропия}
Рассмотрим бромиды щелочных металлов LiBr, NaBr, KBr, RbBr, CsBr. Первые
четыре кристаллизуются в структуре NaCl, а CsBr~— в структуре CsCl.

\begin{opredelenie}
\term{Изоморфизм}~— сохранение типа структуры при замене одних
структурных частиц другими, химически родственными (LiBr\,$\to$\,NaBr\,$\to$\,KBr\,$\to$\,RbBr).
\term{Морфотропия}~— скачкообразное изменение типа структуры в ряду
таких замен (RbBr\,$\to$\,CsBr).
\end{opredelenie}

Геометрическая причина морфотропии видна из табл.~\ref{tab:bromides} и
рис.~\ref{fig:rr-scale}. У LiBr, NaBr, KBr отношение радиусов попадает в
октаэдрический интервал $0{,}414$--$0{,}732$ (у LiBr~— на самой нижней
границе). У RbBr оно ($0{,}76$) уже чуть выше границы, но структура
NaCl ещё сохраняется~— соединение «на пределе» и, как мы увидим, легко
перестраивается. Только у CsBr отношение $0{,}84$ явно превышает
$0{,}732$, и выгоднее становится кубическая координация.

\begin{table}[ht]
\centering
\caption{Геометрическая основа изоморфизма и морфотропии в бромидах
щелочных металлов ($r_{\text{Br}^-}=0{,}196$~нм)}
\label{tab:bromides}
\small
\begin{tabular}{@{}lccccc@{}}
\toprule
Катион & Li$^+$ & Na$^+$ & K$^+$ & Rb$^+$ & Cs$^+$ \\
\midrule
$r_\text{к}$, нм & 0,078 & 0,098 & 0,133 & 0,149 & 0,165 \\
$r_\text{к}/r_\text{а}$ & 0,40 & 0,50 & 0,68 & 0,76 & 0,84 \\
структура & NaCl & NaCl & NaCl & NaCl & CsCl \\
\bottomrule
\end{tabular}
\end{table}

\begin{figure}[ht]
\centering
\begin{tikzpicture}[x=13cm, y=1cm]
  \fill[fmGray!15] (0.155,0) rectangle (0.225,0.7);
  \fill[fmGreen!18] (0.225,0) rectangle (0.414,0.7);
  \fill[fmBlue!15] (0.414,0) rectangle (0.732,0.7);
  \fill[fmAmber!20] (0.732,0) rectangle (1,0.7);
  \node[font=\scriptsize] at (0.19,0.35) {3};
  \node[font=\scriptsize] at (0.32,0.35) {4: тетраэдр};
  \node[font=\scriptsize] at (0.573,0.35) {6: октаэдр};
  \node[font=\scriptsize] at (0.866,0.35) {8: куб};
  \draw[line width=0.5pt] (0.1,0) -- (1,0);
  \foreach \t in {0.155,0.225,0.414,0.732,1}{\draw (\t,0) -- (\t,-0.12) node[below, font=\scriptsize]{\t};}
  % бромиды — над шкалой; цвет = фактическая координация
  \foreach \v/\n/\c in {0.40/LiBr/fmBlue, 0.50/NaBr/fmBlue, 0.68/KBr/fmBlue, 0.76/RbBr/fmBlue, 0.84/CsBr/fmAmber}{
    \draw[\c, line width=0.6pt] (\v,0.7) -- (\v,1.05);
    \fill[\c] (\v,1.05) circle (1.3pt);
    \node[font=\scriptsize, above] at (\v,1.1) {\n};}
  % другие соединения — под шкалой
  \foreach \v/\n/\c in {0.46/ZnS/fmGreen, 0.54/NaCl/fmBlue, 0.78/CaF$_2$/fmAmber, 0.91/CsCl/fmAmber}{
    \draw[\c, line width=0.6pt] (\v,-0.5) -- (\v,-0.8);
    \fill[\c] (\v,-0.8) circle (1.3pt);
    \node[font=\scriptsize, below] at (\v,-0.85) {\n};}
  \node[font=\scriptsize, anchor=east, align=right] at (0.1,0.35) {к.\,ч.};
\end{tikzpicture}
\caption{Шкала отношения радиусов $r_\text{к}/r_\text{а}$ с интервалами
устойчивости координаций. Цвет точки~— фактическая координация катиона
(синий~— 6, жёлтый~— 8, зелёный~— 4). RbBr и CsBr~— пример морфотропии; ZnS~—
исключение из правила}
\label{fig:rr-scale}
\end{figure}

\paragraph{Полиморфизм ионных кристаллов}
Соединения с отношением радиусов вблизи границы легко меняют структуру
при изменении внешних условий.
\begin{itemize}
\item \emph{Под давлением.} RbBr при обычных условиях имеет структуру
NaCl, но при высоком давлении переходит в структуру CsCl: катионы
перемещаются из октаэдрических пустот в кубические. Координационное число
растёт с 6 до 8, упаковка становится плотнее, и объём скачком
уменьшается~— примерно на 15\,\%. Давление «выигрывает» на уменьшении
объёма. Так же ведут себя RbCl и RbI.
\item \emph{При нагреве.} Сам CsCl выше $469$\,°C переходит в структуру
типа NaCl: катионы Cs$^+$ занимают октаэдрические, а не кубические
пустоты\ispr{В оригинале (с.~367) эта модификация названа
низкотемпературной. На самом деле структура типа NaCl у CsCl устойчива
при \emph{высоких} температурах~— выше $469$\,°C и до плавления.}.
Причина: с температурой радиусы разных ионов меняются по-разному, и
отношение $r_\text{к}/r_\text{а}$, бывшее близким к критическому, может
перейти через границу.
\end{itemize}

% ----------------------------------------------------------------------
\subsubsection{Силы Ван-дер-Ваальса. Молекулярные кристаллы}\label{sssec:vdw}

Силами Ван-дер-Ваальса связаны кристаллы благородных газов и кристаллы
из насыщенных молекул (H$_2$, N$_2$, O$_2$, CH$_4$, органические
молекулы)~— внутри таких молекул атомы связаны ковалентно, а между
молекулами действуют только ван-дер-ваальсовы силы.

Это \vazhno{самая слабая} из всех связей. Энергия диссоциации молекулы
H$_2$ около $4{,}5$~эВ, а чтобы разрушить кристалл водорода (испарить
его), достаточно теплоты сублимации около $0{,}01$~эВ на молекулу~— в
сотни раз меньше\ispr{В оригинале (с.~368)~— «в 200 раз». Теплота
сублимации твёрдого водорода около 1~кДж/моль, т.\,е. $\approx0{,}01$~эВ
на молекулу, что примерно в 400 раз меньше энергии диссоциации.}. Тем
не менее именно эти силы определяют все свойства молекулярных
кристаллов.

\paragraph{Происхождение}
Нейтральные атомы и молекулы поляризуются: в электрическом поле их
электронное облако смещается относительно ядер, и появляется
\emph{электрический дипольный момент}. Ван-дер-ваальсово притяжение~—
это электростатическое взаимодействие таких диполей. Различают три
эффекта (рис.~\ref{fig:vdw}).
\begin{enumerate}
\item \term{Ориентационный} (взаимодействие постоянных диполей). У
многих ковалентных молекул есть \emph{постоянный} дипольный момент: атомы
различаются электроотрицательностью, а сама молекула несимметрична
(NH$_3$, H$_2$O, HCl). Если диполи выстроятся так, что «минус» одного
находится рядом с «плюсом» другого, молекулы притягиваются, и энергия
такого упорядоченного расположения заметно ниже, чем хаотического.
\item \term{Индукционный} (дополнение к первому). Диполь создаёт поле,
которое \emph{наводит} дополнительный дипольный момент в соседней
молекуле (электронное облако подвижно). Поэтому дипольный момент молекулы
в кристалле больше, чем у изолированной молекулы.
\item \term{Дисперсионный} (взаимодействие мгновенных диполей). У атомов
благородных газов электронные облака сферически симметричны, постоянного
диполя нет. Но симметрично облако лишь \emph{в среднем}: в каждый момент
электроны расположены несимметрично, и атом имеет \emph{мгновенный}
дипольный момент. Если мгновенные диполи соседних атомов меняются
согласованно (синхронно), атомы притягиваются.
\end{enumerate}

\begin{figure}[ht]
\centering
\begin{tikzpicture}[scale=1.2]
  % а — ориентационный
  \begin{scope}
    \vdwDip{(0,0.45)}{0}\vdwDip{(1.45,0.45)}{0}
    \vdwDip{(0,-0.45)}{180}\vdwDip{(1.45,-0.45)}{180}
    \node[font=\small] at (0.72,-1.45) {\textit{а} ориентационный};
  \end{scope}
  % б — индукционный
  \begin{scope}[shift={(4.3,0)}]
    \vdwDip{(0,0)}{0}
    \draw[atomLine, line width=0.5pt, dashed] (1.55,0) circle (0.42);
    \vdwDip[0.8]{(1.5,0)}{0}
    \fill (1.55,0) circle (0.9pt);
    \draw[-{Stealth[length=1.6mm]}, fmGray] (1.55,0.62) -- (1.3,0.62);
    \node[font=\scriptsize, fmGray, align=center] at (1.55,1.05) {наведённый диполь};
    \node[font=\small] at (0.75,-1.45) {\textit{б} индукционный};
  \end{scope}
  % в — дисперсионный
  \begin{scope}[shift={(8.6,0)}]
    \foreach \y/\ang/\t in {0.45/0/$t_1$, -0.45/180/$t_2$}{
      \vdwDip[0.85]{(0,\y)}{\ang}\vdwDip[0.85]{(1.35,\y)}{\ang}
      \node[font=\scriptsize, fmGray] at (2.25,\y) {\t};}
    \node[font=\small] at (0.75,-1.45) {\textit{в} дисперсионный};
  \end{scope}
\end{tikzpicture}
\caption{Три составляющие сил Ван-дер-Ваальса:
\textit{а}~— притяжение постоянных диполей, выстроенных «плюс к минусу»;
\textit{б}~— диполь наводит диполь в соседней молекуле (пунктир~—
исходное симметричное облако); \textit{в}~— мгновенные диполи атомов
меняются синхронно (показаны два момента времени $t_1$ и $t_2$)}
\label{fig:vdw}
\end{figure}

\begin{podrobnee}[почему дисперсионная энергия $\propto-1/r^6$]
Мгновенный диполь $p_1$ первого атома создаёт на расстоянии $r$ поле
$\mathcal E\propto p_1/r^3$. Это поле наводит во втором атоме диполь
$p_2=\alpha\mathcal E\propto\alpha p_1/r^3$, где $\alpha$~—
поляризуемость. Энергия наведённого диполя в поле
$U\propto-p_2\mathcal E\propto-\alpha p_1^2/r^6$. Среднее $\overline{p_1}=0$,
но $\overline{p_1^2}\neq0$, поэтому притяжение не исчезает при
усреднении. Отсюда показатель $m=6$ в потенциале \eqref{eq:pair-pot}.
\end{podrobnee}

Дисперсионные силы \vazhno{универсальны}: они действуют между любыми
атомами, ионами и молекулами. При сильной химической связи они~—
небольшая добавка, которой часто пренебрегают; в молекулярных кристаллах
составляют значительную долю межмолекулярного взаимодействия; а в
кристаллах благородных газов~— это \emph{единственная} сила притяжения.

\paragraph{Прочность и структуры}
Мерой прочности связи могут служить температуры плавления и кипения
(табл.~\ref{tab:noble}). У благородных газов они очень низки. От неона к
ксенону они заметно растут: растёт число электронов в оболочке, а с ним
поляризуемость атома и дисперсионное притяжение.

\begin{table}[ht]
\centering
\caption{Температуры плавления и кипения благородных газов}
\label{tab:noble}
\small
\begin{tabular}{@{}lcccc@{}}
\toprule
 & Ne & Ar & Kr & Xe \\
\midrule
число электронов & 10 & 18 & 36 & 54 \\
$\Tpl$, °C & $-248{,}7$ & $-189{,}3$ & $-157$ & $-111{,}8$ \\
$T_\text{кип}$, °C & $-246{,}3$ & $-185{,}6$ & $-153$ & $-108{,}0$ \\
\bottomrule
\end{tabular}
\end{table}

Ван-дер-ваальсова связь, как металлическая и ионная, \vazhno{не имеет
направленности и насыщаемости}. Энергия тем ниже, чем ближе частицы, и
молекулярные кристаллы образуют плотнейшие упаковки: неон и аргон~—
ГЦК, молекулы азота N$_2$~— ГПУ, молекулы метана CH$_4$~— ГЦК.

% ----------------------------------------------------------------------
\subsubsection{Сравнение типов связи}\label{sssec:bond-compare}

Сведём всё вместе (табл.~\ref{tab:bond-compare}). Энергии связи даны в
расчёте на атом~— это энергия, нужная, чтобы разобрать кристалл на
нейтральные атомы.

\begin{table}[ht]
\centering
\caption{Четыре типа связи в кристаллах}
\label{tab:bond-compare}
\small
\begin{tabularx}{\textwidth}{@{}>{\raggedright\arraybackslash}p{22mm}*{4}{>{\raggedright\arraybackslash}X}@{}}
\toprule
 & Ковалентная & Металлическая & Ионная & Ван-дер-ваальсова \\
\midrule
Механизм & общая пара электронов с противоположными спинами & ионные остовы в «газе» обобществлённых электронов & переход электрона и кулоновское притяжение ионов & притяжение постоянных, наведённых и мгновенных диполей \\
\addlinespace
Направлен\-ность, насыщае\-мость & есть & нет & нет & нет \\
\addlinespace
Структуры (к.\,ч.) & по правилу $8-N$; алмаз (4) & ГЦК, ГПУ (12), ОЦК (8) & анионы~— упаковка, катионы~— в пустотах (4, 6, 8) & ГЦК, ГПУ (12) \\
\addlinespace
$E_\text{св}$, эВ/атом & C (алмаз) 7,4; Si~4,6; Ge~3,9 & Na~1,1; Fe~4,3; W~8,9 & NaCl~3,3; LiF~4,4 & Ne~0,02; Ar~0,08; Xe~0,16 \\
\bottomrule
\end{tabularx}
\end{table}

\begin{itogi}
\begin{itemize}
\item Все типы связи имеют электрическую природу; энергия пары атомов~—
сумма притяжения и отталкивания, $U=-A/r^m+B/r^n$; минимум задаёт
равновесное расстояние и энергию связи.
\item \emph{Ковалентная} связь~— общая электронная пара; связь
направленна и насыщаема, отсюда рыхлые структуры по правилу $8-N$
(алмаз: к.\,ч.~4).
\item \emph{Металлическая}~— ионные остовы в электронном газе;
ненаправленна и ненасыщаема, отсюда плотные упаковки ГЦК, ГПУ и близкая к
ним ОЦК.
\item \emph{Ионная}~— переход электрона и кулоновское притяжение; анионы
образуют упаковку, катионы занимают пустоты. Тип пустот подсказывает
отношение $r_\text{к}/r_\text{а}$: $0{,}225$~— $0{,}414$~— $0{,}732$~—
границы тетраэдрической, октаэдрической и кубической координаций.
\item \emph{Ван-дер-ваальсова}~— притяжение диполей ($\propto-1/r^6$),
самая слабая (сотые и десятые доли эВ на атом); плотнейшие упаковки.
\item Полиморфизм и морфотропия возникают, когда условия (температура,
давление, замена атомов) переводят вещество через границу выгодности
структур.
\end{itemize}
\end{itogi}

\begin{proverka}
\item Электроотрицательности по Мулликену у натрия и хлора~— около $2{,}8$
и $8{,}3$~эВ. Какой атом отдаст электрон и почему?
\item Почему правило $D_{AB}\approx(D_{AA}+D_{BB})/2$ «не работает» для NaCl
и что из этого следует?
\item Выведите радиус тетраэдрической пустоты $\sqrt{3/2}-1$.
\item Ионный радиус Li$^+$ равен $0{,}078$~нм, I$^-$~— $0{,}220$~нм. Какую
координацию лития можно ожидать в LiI?
\item Почему RbBr под давлением переходит в структуру CsCl, а не
наоборот?
\item Почему температура плавления растёт в ряду Ne\,--\,Ar\,--\,Kr\,--\,Xe?
\end{proverka}
